\documentclass[aspectratio=169, 14pt]{beamer}
\usepackage{stampfly_slides}

\lesson{0}{環境セットアップ}{Environment Setup}

\begin{document}

% ------------------------------------------------------------------
\begin{frame}
  \titlepage
\end{frame}

% ------------------------------------------------------------------
\begin{frame}{ワークショップカリキュラム / Workshop Curriculum}
  \begin{table}
    \centering
    \begin{tabular}{lll}
      \hline
      \textbf{Day} & \textbf{Lesson} & \textbf{テーマ} \\
      \hline
      Day 1 & 0\,--\,4   & 環境セットアップ + ハードウェア入門 \\
      Day 2 & 5\,--\,7   & フィードバック制御 + 初フライト \\
      Day 3 & 8\,--\,9   & モデリング + 姿勢推定 \\
      Day 4 & 10\,--\,11 & Python SDK + 競技会 \\
      \hline
    \end{tabular}
  \end{table}
\end{frame}

% ------------------------------------------------------------------
\begin{frame}{StampFly ハードウェア / Hardware}
  \begin{table}
    \centering
    \begin{tabular}{ll}
      \hline
      \textbf{項目} & \textbf{仕様} \\
      \hline
      MCU    & ESP32-S3 (M5Stamp S3) \\
      IMU    & BMI270 (6軸 400\,Hz) \\
      気圧   & BMP280 \\
      ToF    & VL53L3CX \\
      質量   & 35\,g（バッテリ含む） \\
      通信   & ESP-NOW + WiFi \\
      バッテリ & LiPo 1S 3.7\,V \\
      \hline
    \end{tabular}
  \end{table}
\end{frame}

% ------------------------------------------------------------------
\begin{frame}{エコシステム全体像 / Ecosystem Overview}
  \begin{columns}[T]
    \begin{column}{0.52\textwidth}
      \centering
      \vspace{-2mm}
      \resizebox{0.95\columnwidth}{!}{\documentclass[border=10pt]{standalone}
\usepackage{tikz}
\usetikzlibrary{arrows.meta, positioning, shapes.geometric, fit}
\usepackage{luatexja}
\usepackage[match]{luatexja-fontspec}
% Cross-platform font: Hiragino (macOS) → Yu Gothic (Windows) → Noto (Linux)
\usepackage{ifthen}
\newboolean{jfontset}\setboolean{jfontset}{false}
\IfFontExistsTF{Hiragino Kaku Gothic ProN}{%
  \setmainjfont{Hiragino Kaku Gothic ProN}%
  \setsansjfont{Hiragino Kaku Gothic ProN}%
  \setboolean{jfontset}{true}}{}
\ifthenelse{\boolean{jfontset}}{}{%
  \IfFontExistsTF{Yu Gothic}{%
    \setmainjfont{Yu Gothic}%
    \setsansjfont{Yu Gothic}%
    \setboolean{jfontset}{true}}{}}
\ifthenelse{\boolean{jfontset}}{}{%
  \IfFontExistsTF{Noto Sans CJK JP}{%
    \setmainjfont{Noto Sans CJK JP}%
    \setsansjfont{Noto Sans CJK JP}}{}}

\begin{document}

\definecolor{layer1}{RGB}{0, 100, 180}   % firmware
\definecolor{layer2}{RGB}{40, 140, 60}   % control
\definecolor{layer3}{RGB}{180, 120, 0}   % analysis
\definecolor{layer4}{RGB}{0, 120, 200}   % tools
\definecolor{layer5}{RGB}{120, 80, 160}  % simulator
\definecolor{layer6}{RGB}{180, 60, 40}   % protocol

\begin{tikzpicture}[
    layer/.style={
      rectangle, rounded corners=4pt, draw=#1, line width=1.5pt,
      fill=#1!10, minimum width=75mm, minimum height=14mm,
      font=\bfseries\large, text=#1, align=center
    },
    label/.style={
      font=\small, text=black!70, align=left, anchor=west
    },
  ]

  % Layers (bottom to top)
  \node[layer=layer6] (L6) at (0, 0)    {\texttt{protocol/}\hspace{1em}通信仕様};
  \node[layer=layer1, above=5mm of L6]  (L1) {\texttt{firmware/}\hspace{1em}組込みコード};
  \node[layer=layer2, above=5mm of L1]  (L2) {\texttt{control/}\hspace{1em}制御設計};
  \node[layer=layer5, above=5mm of L2]  (L5) {\texttt{simulator/}\hspace{1em}シミュレータ};
  \node[layer=layer3, above=5mm of L5]  (L3) {\texttt{analysis/}\hspace{1em}データ解析};
  \node[layer=layer4, above=5mm of L3]  (L4) {\texttt{tools/}\hspace{1em}開発ツール};

  % Right-side descriptions
  \node[label, right=12mm of L6] {SSOT: YAML → コード自動生成};
  \node[label, right=12mm of L1] {ESP32-S3 / ESP-IDF / C++};
  \node[label, right=12mm of L2] {PID, MPC, SIL};
  \node[label, right=12mm of L5] {Genesis 3D 物理};
  \node[label, right=12mm of L3] {Jupyter / Python};
  \node[label, right=12mm of L4] {sf CLI};

\end{tikzpicture}
\end{document}
}
    \end{column}
    \begin{column}{0.45\textwidth}
      \footnotesize
      \begin{itemize}\setlength{\itemsep}{2pt}
        \item \texttt{firmware/} --- 組込みコード
        \item \texttt{control/} --- 制御設計
        \item \texttt{simulator/} --- 3D シミュレータ
        \item \texttt{analysis/} --- データ解析
        \item \texttt{tools/} --- sf CLI
        \item \texttt{protocol/} --- 通信仕様
      \end{itemize}
    \end{column}
  \end{columns}
\end{frame}

% ------------------------------------------------------------------
\begin{frame}{開発ツール sf CLI / Development Tool}
  \begin{table}
    \centering\small
    \begin{tabular}{lll}
      \hline
      \textbf{カテゴリ} & \textbf{コマンド例} & \textbf{説明} \\
      \hline
      ビルド         & \code{sf build} / \code{sf flash}    & ファームウェア開発 \\
      診断           & \code{sf doctor} / \code{sf monitor}  & デバッグ \\
      ログ           & \code{sf log wifi} / \code{sf log viz} & テレメトリ取得・可視化 \\
      シミュレータ   & \code{sf sim run}                     & 仮想環境で練習 \\
      キャリブレーション & \code{sf cal gyro/accel/mag}      & センサ校正 \\
      \hline
    \end{tabular}
  \end{table}

  \vspace{0.5em}

  \begin{block}{ワークフロー}
    \code{sf lesson switch N} → \code{sf build workshop} → \code{sf flash workshop -m}
  \end{block}
\end{frame}

% ------------------------------------------------------------------
\begin{frame}{今日のゴール / Today's Goal}
  \begin{block}{目標}
    ワークショップファームウェアをビルド・書き込み・動作確認する
  \end{block}

  \vspace{1em}

  \begin{enumerate}
    \item \code{sf build workshop} でビルド
    \item \code{sf flash workshop -m} で書き込み
    \item シリアルモニタで \texttt{"Hello StampFly!"} を確認
  \end{enumerate}
\end{frame}

% ------------------------------------------------------------------
\begin{frame}{開発フロー / Development Flow}
  \begin{center}
    \resizebox{0.85\textwidth}{!}{\documentclass[border=10pt]{standalone}
\usepackage{tikz}
\usetikzlibrary{arrows.meta, positioning, shapes.geometric}

\definecolor{blue1}{RGB}{0, 100, 180}
\definecolor{green1}{RGB}{50, 140, 50}
\definecolor{yellow1}{RGB}{190, 150, 0}
\definecolor{purple1}{RGB}{120, 80, 150}
\definecolor{gray1}{RGB}{120, 120, 120}

\begin{document}
\begin{tikzpicture}[
    step/.style={
      rectangle, rounded corners=5pt, draw=#1, line width=1.5pt,
      fill=#1!8, minimum width=42mm, minimum height=18mm,
      font=\bfseries\large, text=#1, align=center
    },
    desc/.style={font=\small, text=gray1, align=center, below=2mm},
    arr/.style={-{Stealth[length=4mm]}, line width=2pt, #1},
    feedback/.style={-{Stealth[length=3mm]}, line width=1.2pt, gray1, dashed},
  ]

  % Steps
  \node[step=blue1] (s1) at (0, 0) {sf lesson\\switch N};
  \node[step=green1, right=22mm of s1] (s2) {sf build\\workshop};
  \node[step=yellow1, right=22mm of s2] (s3) {sf flash\\workshop -m};
  \node[step=purple1, right=22mm of s3] (s4) {Serial\\Monitor};

  % Arrows
  \draw[arr=blue1] (s1) -- (s2);
  \draw[arr=green1] (s2) -- (s3);
  \draw[arr=yellow1] (s3) -- (s4);

  % Descriptions
  \node[desc] at (s1.south) {student.cpp $\to$\\user\_code.cpp};
  \node[desc] at (s2.south) {ESP-IDF\\compile};
  \node[desc] at (s3.south) {USB flash +\\monitor};
  \node[desc] at (s4.south) {verify\\output};

  % Feedback loop
  \draw[feedback] (s4.south) -- ++(0, -1.5) -| node[pos=0.25, below, font=\small\itshape, text=gray1] {modify code \& rebuild} (s1.south west);

  % Title
  \node[font=\LARGE\bfseries, above=15mm of s2.north east] {Build $\to$ Flash $\to$ Monitor Flow};

\end{tikzpicture}
\end{document}
}
  \end{center}
\end{frame}

% ------------------------------------------------------------------
\begin{frame}{sf CLI とは / What is sf CLI?}
  \begin{itemize}
    \item StampFly 開発専用のコマンドラインツール
    \item ESP-IDF のコマンドをシンプルに実行できる
  \end{itemize}

  \vspace{0.5em}

  \begin{table}
    \centering\small
    \begin{tabular}{ll}
      \hline
      \textbf{コマンド} & \textbf{説明} \\
      \hline
      \code{sf doctor}          & 環境診断 \\
      \code{sf build workshop}  & ワークショップビルド \\
      \code{sf flash workshop -m} & 書き込み＋モニタ \\
      \code{sf lesson switch N} & レッスン切り替え \\
      \code{sf monitor}         & シリアルモニタ \\
      \hline
    \end{tabular}
  \end{table}
\end{frame}

% ------------------------------------------------------------------
\begin{frame}{ワークショップの構造 / Workshop Structure}
  \begin{block}{学生が実装する関数は2つだけ！}
    \begin{itemize}
      \item \api{setup()} --- 起動時に1回呼ばれる
      \item \api{loop\_400Hz(dt)} --- 400Hz（2.5ms毎）で呼ばれる
    \end{itemize}
  \end{block}

  \vspace{0.5em}

  \begin{itemize}
    \item ハードウェアの複雑さは \api{ws::} 名前空間で隠蔽
    \item FreeRTOS, SPI/I2C, センサーフュージョンを意識する必要なし
  \end{itemize}
\end{frame}

% ------------------------------------------------------------------
\begin{frame}[fragile]{実習: Hello StampFly / Hands-on}
  \begin{lstlisting}[style=sfcpp, title={student.cpp}]
#include "workshop_api.hpp"

void setup()
{
    ws::print("Hello StampFly!");
}

void loop_400Hz(float dt)
{
    // Nothing to do in Lesson 0
}
  \end{lstlisting}
  \vspace{-0.8em}
  \begin{exampleblock}{手順}\small
    \textbf{1.}~\code{sf lesson switch 0}\\
    \textbf{2.}~上のコードを \code{student.cpp} に書く\\
    \textbf{3.}~\code{sf build workshop \&\& sf flash workshop -m}
  \end{exampleblock}
\end{frame}

% ------------------------------------------------------------------
\begin{frame}{チェックポイント / Checkpoint}
  \begin{alertblock}{確認事項}
    \begin{itemize}
      \item[$\square$] \code{sf doctor} がエラーなしで通る
      \item[$\square$] ビルドが成功する
      \item[$\square$] シリアルモニタに \texttt{"Hello StampFly!"} が表示される
    \end{itemize}
  \end{alertblock}

  \vspace{1em}

  \begin{block}{次のレッスン}
    Lesson 1: モータ制御 --- モータを回してプロペラの動きを確認
  \end{block}
\end{frame}

\end{document}
